% Conclusion - Reframed: Confident, not defensive

We discover that KG uncertainty methods are systematically overconfident on zero-evidence queries: among top-confident predictions, zero-evidence queries appear at 2--2.5$\times$ the base rate. This affects 11--25\% of test queries across seven benchmarks---a substantial failure mode hidden by aggregate evaluation.

The cause is structural: entity-level uncertainty cannot detect relation-specific novelty. High-frequency entities receive low uncertainty regardless of which relation is queried. Theorem~\ref{thm:impossibility} formalizes this, and our experiments confirm it holds across variational embeddings, ensembles, energy scoring, and foundation models (ULTRA achieves 0.29 AUROC---below random).

The fix---tracking entity-relation coverage---is obvious in hindsight. Yet no major uncertainty method implements it: UKGE, BEUrRE, GP-KGE, Energy scoring, Deep Ensembles, and ULTRA all lack this mechanism. The contribution is not inventing a complex solution but exposing that a simple one was universally missing.

\paragraph{Practical recommendations.}
\begin{enumerate}
\item \textbf{Track entity-relation coverage.} A hash table catches the failure mode invisible to learned methods. For billion-entity KGs, Bloom filters reduce storage 5$\times$ with $<$0.2pp recall loss.
\item \textbf{Flag zero-coverage queries.} When $c(e,r)=0$, the model has no evidence for this combination.
\item \textbf{Report stratified metrics.} Aggregate AUROC hides category-specific failures.
\end{enumerate}

\paragraph{Limitations.}
Coverage tracking fails on role-shift OOD (GDELT: 0.57 AUROC) where entities use atypical relations but have full coverage. This requires relation-conditioned methods beyond our scope. The approach assumes transductive settings; inductive KGs need different mechanisms.

\paragraph{Reproducibility.}
Code, models, and data splits released upon publication.
